\section{Bible Conference at Eagle Lake}

“Folkebladet,” May 01, 1918

The Bible Conference at Eagle Lake was held, according to the announcement, from March 19–21 and was very successful. It would, however, have benefited the cause of God’s kingdom if more people had come out to such blessed meetings.

The opening sermon was preached by K. T. Nytten. Pastor Gynild spoke on the discussion topic, since the one appointed to introduce it had been unavoidably prevented from attending. With regard to the lectures that were given, it need only be said that they will be printed, since the assembly considered them worthy of it. Pastor Ludvig Pedersen gave a lecture on “The Congregation’s Origin, Composition, and Equipment.”

The next day the discussion topic, Ephesians 4:1–16, was introduced by C. M. Hanson. Prof. C. P. Harbo spoke on the subject. In the afternoon Pastor C. C. Gynild gave a lecture on “The Responsibility of the Individual Congregation in Carrying Out the Common Tasks.” He laid emphasis on the fact that the common tasks must be carried out in freedom; but it costs a struggle to win freedom and a struggle to preserve the freedom that has been won; and not least does it cost a struggle to preserve faith in freedom, since many are very much afraid of it.

Prof. C. P. Harbo then began his lecture on “The Distinctive Position of the Lutheran Free Church, Its Calling and Task.”

In the evening Pastor A. Olson gave a lecture on “Joseph.” It was planned especially as a lecture for the young people.

The next morning a business meeting was held, since, as will be remembered, this meeting takes the place of the district’s winter meeting. The minutes of the last meeting were read and adopted. The treasurer reported. It was resolved that the next district meeting be held in Eidskog congregation, Pastor K. Knutson’s charge, from June 23–25. It was resolved that the youth convention be held in Kongsvinger congregation, Pastor Olson’s charge, from June 28–30.

Pastor E. O. Larson will introduce the discussion topic at the next district meeting.

The music committee for the youth convention, which was elected last year, will continue this year also, with the change that Pastor L. Pedersen is elected in place of Mrs. Pastor Konsterlie, who is now in China. Pastor E. M. Hanson will introduce the topic at the youth convention. Pastor L. Pedersen will preach the opening sermon. The meeting was adjourned.

Prof. C. P. Harbo then continued his lecture, demonstrating from Scripture and history how the Lutheran Free Church had come to occupy this distinctive position. He laid emphasis on the memorable words which Prof. Sverdrup Sr. spoke on that memorable evening of October 10, 1894, in Trinity Church: “Therefore Augsburg was built, so that the Norwegian awakening might have a way opened to it and carry on its work in the new land.”

Pastor K. Knutson, chairman of the district, gave a lecture in the afternoon on “The Signs of the Times.” But, as has been said, we will not anticipate matters by telling what each one said, since the lectures will be printed; and the brothers who gave the lectures are hereby requested to put their lectures into a form fit for printing as soon as possible and send them to the undersigned. Finally Pastor E. O. Larson, the pastor of the place, spoke, and so the great meeting came to an end.

Besides those mentioned above, Pastor N. Nilsen and his son, Pastor N. Nilsen, were also present at the meeting.

Before the sessions, a prayer meeting was held each morning, and many were there who took part.

And so, heartfelt thanks to all. Thanks also to the women for the excellent hospitality.

E. M. Hanson,
Secretary